\chapter{Updates}

A release has parts that are updated in different ways:

\begin{center}
\begin{tabularx}{\linewidth}{@{}l l X@{}}
\toprule
\textbf{Part} & \textbf{Where it lives} & \textbf{How it is updated} \\
\midrule
FPGA bitstream & the FPGA's configuration flash & JTAG (FlashPro), with physical access \\
Driver and service & microSD card & over the network, with the installers \\
Client library & the user's PC & \cmd{pip install} of the new wheel \\
\bottomrule
\end{tabularx}
\end{center}

\begin{advertencia}[Board limitation]
PolarFire SoC can update its bitstream in the field (\emph{Auto Update} and
\emph{IAP}), but both need an SPI flash memory connected to the chip's System
Controller, which the Discovery Kit does not have. On this board the
bitstream can only be updated over JTAG. A product board should include that
memory.
\end{advertencia}

\section{Compatibility}

\begin{itemize}
  \item \cmd{/v1/device} reports the versions of the board's three parts;
    \cmd{bitstream} identifies the release and the exact commit.
  \item Within protocol v1, a client of any version works with a device of
    any version: the protocol only grows by addition.
\end{itemize}

\section{File integrity}

Every release publishes a \cmd{SHA256SUMS} file. Before installing:

\begin{lstlisting}[language=shell]
sha256sum -c SHA256SUMS
\end{lstlisting}
